%==============================================================
% BAB III  METODE
%==============================================================

\chapter{METODE}

%--------------------------------------------------------------
\section{Data dan Strategi Pembagian Data}
%--------------------------------------------------------------

\subsection{Dataset Utama: Zhao2024}

Dataset utama yang digunakan dalam penelitian ini adalah koleksi citra fundus bayi yang
diterbitkan oleh \citet{Zhao2024}, selanjutnya disebut Zhao2024. Dataset ini terdiri dari
1.099 citra fundus bayi prematur dalam lima kelas: Normal, Stage 1, Stage 2, Stage 3, dan
bekas luka laser (\textit{laser scars}). Rincian distribusi kelas ditunjukkan pada Tabel
\ref{tab:distribusi_kelas}.

\begin{table}[H]
\centering
\caption{Distribusi kelas pada dataset Zhao2024}
\label{tab:distribusi_kelas}
\begin{tabular}{l c}
\toprule
\textbf{Kelas} & \textbf{Jumlah Citra} \\
\midrule
Normal            & 236  \\
Stage 1           & 94   \\
Stage 2           & 165  \\
Stage 3           & 261  \\
Laser Scars       & 343  \\
\midrule
\textbf{Total}    & \textbf{1.099} \\
\bottomrule
\end{tabular}
\end{table}

Analisis mendalam terhadap dataset menemukan adanya 12 citra yang identik secara byte
(\textit{byte-identical duplicates}) dalam 6 kelompok duplikat. Yang lebih penting, tiga
dari kelompok tersebut merupakan konflik label antar-kelas: citra yang persis sama tercatat
di bawah dua kelas yang berbeda. Pasangan yang teridentifikasi adalah \textit{Stage\_2\_ROP\_154.jpg}
$\equiv$ \textit{Stage\_3\_ROP\_210.jpg}, \textit{Stage\_2\_ROP\_52.jpg} $\equiv$
\textit{Stage\_3\_ROP\_254.jpg}, dan \textit{Stage\_1\_ROP\_42.jpg} $\equiv$
\textit{Stage\_2\_ROP\_4.jpg}. Duplikat ini dibiarkan ada dalam dataset dan ditangani melalui
protokol pengelompokan (\textit{grouping}), bukan dengan menghapus data.

\subsection{Dataset Sekunder: Agrawal2021}

Dataset sekunder yang digunakan adalah bagian HVDROPDB-BV dari dataset yang diterbitkan oleh
\citet{Agrawal2021}. Dataset ini terdiri dari 100 pasang citra fundus bayi beserta masker
pembuluh darah yang dianotasi oleh pakar, yang terbagi atas 50 citra dari kamera RetCam dan
50 citra dari kamera Neo. Dataset ini digunakan semata-mata untuk mengevaluasi kualitas
pipeline segmentasi pembuluh darah secara kuantitatif menggunakan metrik Dice dan tidak
terlibat dalam pelatihan atau pengujian model klasifikasi.

\subsection{Strategi Pembagian Data dan Pengelompokan}

Karena dataset Zhao2024 tidak menyediakan identitas pasien pada nama file, pembagian data
secara langsung berbasis identitas pasien tidak dapat dilakukan. Sebagai solusi, analisis
kemiripan antar citra menggunakan korelasi silang ternormalisasi (\textit{Normalized Cross-
Correlation}, NCC $\geq$ 0,99) dilakukan untuk mendeteksi pasangan citra yang berasal dari
mata yang sama. Proses ini menghasilkan 721 kelompok dari total 1.099 citra yang kemudian
disimpan dalam berkas \textit{clean\_manifest.json}.

Validasi silang 5-\textit{fold} menggunakan \textit{StratifiedGroupKFold} diterapkan agar
setiap kelompok selalu berada sepenuhnya di satu sisi batas pelatihan--pengujian. Strategi ini
memastikan bahwa model benar-benar mempelajari pola penyakit, bukan pola citra dari mata
tertentu.

%--------------------------------------------------------------
\section{Pipeline Segmentasi Pembuluh Darah}
%--------------------------------------------------------------

Pipeline segmentasi pembuluh darah dirancang berdasarkan prinsip detektor berbasis filter
klasik tanpa komponen \textit{deep learning}. Pipeline ini dioptimalkan melalui serangkaian
percobaan terkontrol menggunakan dataset HVDROPDB-BV, dengan evaluasi menggunakan nilai Dice
pada set pengujian yang terpisah (40\% dari 100 citra). Gambar \ref{fig:vessel_segmentation_rounds}
menunjukkan kemajuan nilai Dice sepanjang proses optimasi.

\begin{figure}[H]
  \centering
  \includegraphics[width=0.8\linewidth]{images/vessel_segmentation_rounds}
  \caption{Kemajuan nilai Dice segmentasi pembuluh darah sepanjang proses optimasi. Konfigurasi
  terbaik \textit{g40\_m60\_fine} mencapai Dice = 0,4739}
  \label{fig:vessel_segmentation_rounds}
\end{figure}

\subsection{Konfigurasi Optimal: \textit{g40\_m60\_fine}}

Hasil terbaik dicapai melalui fusi dua detektor dengan formula:
\begin{equation}
  \text{softmap} = 0{,}40 \times \text{gabor\_tophat} + 0{,}60 \times \text{meijering\_fine}
  \label{eq:softmap}
\end{equation}

\noindent Langkah-langkah pemrosesan adalah sebagai berikut:

\begin{enumerate}
    \item \textbf{Praproses kanal hijau}: Kanal hijau diekstrak dari citra RGB. Piksel di
    luar area \textit{field of view} (FOV) diisi menggunakan nilai piksel FOV terdekat untuk
    mencegah artefak pada tahap filtering berikutnya.
    \item \textbf{Normalisasi latar belakang}: Latar belakang diperkirakan menggunakan filter
    Gaussian berskala besar ($\sigma \approx$ 30 piksel). Kanal hijau dibagi dengan perkiraan
    latar belakang tersebut untuk meratakan pencahayaan.
    \item \textbf{CLAHE pertama}: CLAHE dengan \textit{clip limit} 6,0 dan ukuran \textit{tile}
    $16 \times 16$ diterapkan pada kanal hijau yang sudah dinormalisasi.
    \item \textbf{Preprocessing top-hat}: Filter \textit{top-hat} morfologis diterapkan pada
    kanal yang sudah diproses untuk mempertajam struktur pembuluh darah gelap sebelum filtering
    Gabor.
    \item \textbf{Filter Gabor multiskala}: Bank filter Gabor multiskala (sigma: 1,5; 2,5; 3,5;
    5,0 piksel) dan multiorientasi (15 sudut dari $0°$ hingga $165°$) diterapkan. Respons
    maksimum dari seluruh konfigurasi diambil sebagai peta respons Gabor.
    \item \textbf{Penghalusan dan CLAHE kedua}: Median blur $7 \times 7$ diterapkan untuk
    mengurangi \textit{noise}, diikuti CLAHE kedua (\textit{clip limit} 12,0, tile $12 \times 12$).
    Peta gabor\_tophat diperoleh dari tahap ini.
    \item \textbf{Filter Meijering}: Filter Meijering dengan sigma halus (0,8; 1,4; 2,0; 2,8;
    3,6; 4,5 piksel) diterapkan secara terpisah untuk mendeteksi pembuluh darah yang lebih tipis
    yang terlewat oleh filter Gabor.
    \item \textbf{Fusi dan ambang batas}: Softmap akhir dihitung menggunakan persamaan
    \ref{eq:softmap}. Binarisasi menggunakan ambang persentil P0,16 dengan erosi area FOV,
    diikuti pemilihan tiga komponen terhubung terbesar dan penutupan morfologis ($3 \times 3$).
\end{enumerate}

Proses optimasi ini dimulai dari nilai Dice dasar 0,4469 (konfigurasi awal) dan mencapai
nilai Dice 0,4739 pada iterasi ketujuh. Ringkasan kemajuan ditunjukkan pada Tabel
\ref{tab:vessel_progress}.

\begin{table}[H]
\centering
\caption{Ringkasan kemajuan nilai Dice segmentasi pembuluh darah}
\label{tab:vessel_progress}
\begin{tabular}{l l c c}
\toprule
\textbf{Iterasi} & \textbf{Perubahan Utama} & \textbf{Dice (Test)} & \textbf{$\Delta$ vs Baseline} \\
\midrule
Baseline & Gabor + 2$\times$CLAHE + P0,16 & 0,4469 & --- \\
1 & Optimasi ambang batas dan \textit{cleanup} & 0,4505 & +0,0036 \\
2 & Preprocessing top-hat sebelum Gabor & 0,4600 & +0,0131 \\
3 & Penyesuaian rasio blend top-hat & 0,4599 & (stabil) \\
4 & Fusi Sato + Gabor & 0,4655 & +0,0186 \\
5 & Sigma Sato yang lebih halus & 0,4667 & +0,0198 \\
6 & Tiga cara fusi (Meijering mengganti Sato) & --- & (redudansi) \\
\textbf{7} & \textbf{Fusi Gabor + Meijering (terbaik)} & \textbf{0,4739} & \textbf{+0,0270} \\
8 & Konfigurasi tambahan di luar iterasi 7 & (regresi) & --- \\
\bottomrule
\end{tabular}
\end{table}

%--------------------------------------------------------------
\section{Skenario Representasi Input}
%--------------------------------------------------------------

Penelitian ini membandingkan lima skenario representasi input untuk menguji pengaruh
berbagai bentuk pemrosesan citra terhadap performa model klasifikasi. Gambar
\ref{fig:pipeline_diagram} menunjukkan kerangka umum pipeline penelitian.

\begin{figure}[H]
  \centering
  \includegraphics[width=\linewidth]{images/pipeline_diagram}
  \caption{Kerangka \textit{pipeline} penelitian yang menunjukkan lima skenario input dan
  alur pemrosesan menuju model \textit{TinyResNetV2}}
  \label{fig:pipeline_diagram}
\end{figure}

\vspace{2mm}
\noindent\textbf{S1: RGB Mentah (\textit{Raw RGB})}
Citra fundus asli yang hanya diubah ukurannya menjadi $224 \times 224$ piksel dan dinormalisasi
nilai pikselnya ke rentang $[0, 1]$. Skenario ini berfungsi sebagai \textit{baseline} mutlak
tanpa pemrosesan tambahan apapun.

\vspace{2mm}
\noindent\textbf{S2: RGB Enhanced (CLAHE)}
Citra yang sudah diproses menggunakan koreksi pencahayaan berbasis normalisasi Gaussian,
diikuti CLAHE pada kanal kecerahan dalam ruang warna CIE L*a*b*. Skenario ini menguji
apakah peningkatan kontras secara visual berkontribusi pada performa klasifikasi.

\vspace{2mm}
\noindent\textbf{S3: Vessel Mask}
Citra biner satu kanal yang berisi hasil segmentasi pembuluh darah. Direplikasi menjadi tiga
kanal untuk memenuhi spesifikasi input model. Skenario ini menguji apakah informasi bentuk
dan distribusi pembuluh darah saja---tanpa tekstur retina lainnya---sudah cukup informatif.

\vspace{2mm}
\noindent\textbf{S4: Vessel-Guided Enhanced}
Citra S2 yang dikalikan secara piksel dengan peta probabilitas halus (\textit{soft vesselness
map}) turunan dari S3, sehingga area di sekitar pembuluh darah mendapat penekanan lebih kuat.
Skenario ini menguji apakah penekanan konteks vaskular membantu model memfokuskan perhatiannya
pada area yang relevan.

\vspace{2mm}
\noindent\textbf{S5: Masked Softmap (Terbaik)}
Masukan tiga kanal yang dibentuk oleh: (1) \textit{softmap} pembuluh darah dari konfigurasi
\textit{g40\_m60\_fine}, (2) \textit{softmap} area tepi (\textit{ridge softmap}), dan
(3) kanal hijau ternormalisasi CLAHE yang telah di-masker FOV. Semua kanal bersifat kontinu
(bukan biner) dan dimasker oleh FOV. Skenario ini menghasilkan performa terbaik di antara
seluruh skenario yang diuji.

%--------------------------------------------------------------
\section{Baseline Klasik: Fitur Rekayasa Manual}
%--------------------------------------------------------------

Sebelum membangun model jaringan saraf, sebuah \textit{baseline} berbasis fitur rekayasa manual
(\textit{handcrafted features}) dibangun untuk menetapkan batas bawah performa yang harus
dilampaui. Baseline ini menggunakan 48 fitur yang diekstrak dari setiap citra, mencakup:

\begin{itemize}
    \item \textbf{Fitur morfologi pembuluh darah}: Kerapatan pembuluh darah (\textit{vessel
    density}), proksi tortuositas, anisotropi pembuluh, ukuran dan standar deviasi komponen.
    \item \textbf{Fitur tekstur}: Matriks kookurensi tingkat keabuan (\textit{Gray-Level
    Co-occurrence Matrix}, GLCM) \citep{Khandouzi2022} yang menghasilkan fitur kontras,
    korelasi, energi, dan homogenitas.
    \item \textbf{Fitur statistik intensitas}: Rerata dan standar deviasi intensitas pada
    setiap kanal warna.
    \item \textbf{Fitur area tepi (\textit{ridge})}: Luas, panjang, kerapatan, dan persentil
    respons dari komponen \textit{ridge} yang terdeteksi.
\end{itemize}

Fitur-fitur tersebut kemudian dimasukkan ke dalam tiga jenis model klasifikasi: \textit{Support
Vector Machine} (SVM) dengan kernel RBF, \textit{Random Forest} \citep{Breiman2001} dengan 400
pohon keputusan (\textit{n\_estimators=400, max\_depth=None}), dan \textit{Gradient Boosting}.
Evaluasi dilakukan menggunakan validasi silang 5-\textit{fold} yang distratifikasi berdasarkan kelas.

%--------------------------------------------------------------
\section{Arsitektur Model Klasifikasi: TinyResNetV2}
%--------------------------------------------------------------

Model klasifikasi utama adalah \textit{TinyResNetV2}, sebuah jaringan saraf residual kecil yang
dirancang dari awal (\textit{from scratch}) tanpa menggunakan bobot pra-latih. Arsitektur ini
dipilih karena ukurannya yang proporsional dengan skala dataset sehingga risiko \textit{overfitting}
dapat dikendalikan.

Arsitektur \textit{TinyResNetV2} terdiri dari:
\begin{enumerate}
    \item Lapisan konvolusi awal $3 \times 3$ diikuti \textit{batch normalization} dan fungsi
    aktivasi ReLU.
    \item Tiga tahap \textit{residual block} bertingkat dengan lebar fitur masing-masing 48,
    96, dan 192 kanal. Setiap tahap menerapkan \textit{downsampling} spasial melalui konvolusi
    dengan \textit{stride} 2.
    \item \textit{Global Average Pooling} untuk meringkas peta fitur dua dimensi menjadi vektor
    satu dimensi.
    \item Lapisan \textit{Dropout} untuk regulasi.
    \item Lapisan \textit{Dense} dengan fungsi aktivasi \textit{softmax} untuk lima kelas.
\end{enumerate}

Struktur \textit{residual block} mengikuti pola dasarnya: \textit{conv-BN-ReLU-conv-BN} dengan
hubungan pintas (\textit{skip connection}) yang menjumlahkan masukan blok dengan keluarannya.
Arsitektur ini ditunjukkan pada Gambar \ref{fig:tinyresnet_arch}.

\begin{figure}[H]
  \centering
  \includegraphics[width=0.9\linewidth]{images/tinyresnet_arch}
  \caption{Arsitektur \textit{TinyResNetV2} dengan tiga tahap \textit{residual block} dan lebar
  kanal bertingkat (48/96/192). Hubungan pintas (\textit{skip connection}) memungkinkan gradien
  mengalir langsung melewati setiap blok}
  \label{fig:tinyresnet_arch}
\end{figure}

%--------------------------------------------------------------
\section{Protokol Pelatihan}
%--------------------------------------------------------------

\subsection{Fungsi Kerugian dan Pembobotan Kelas}

Pelatihan menggunakan \textit{focal loss} \citep{Lin2017} dengan parameter $\gamma = 1{,}0$
dan bobot kelas berbanding terbalik dengan frekuensi masing-masing kelas
(\textit{inverse-frequency class weights}). \textit{Focal loss} secara otomatis mengurangi
kontribusi sampel yang mudah diklasifikasikan, sehingga pelatihan lebih berfokus pada sampel
kelas minoritas yang sulit.

\subsection{Optimisasi dan Penjadwalan Laju Belajar}

Optimisasi menggunakan algoritma AdamW \citep{Loshchilov2019} dengan laju belajar awal
$10^{-3}$ dan peluruhan bobot (\textit{weight decay}) $5 \times 10^{-4}$. Laju belajar
dijadwalkan menggunakan strategi \textit{cosine annealing} selama 160 \textit{epoch}.

\subsection{Augmentasi Data}

Augmentasi data diterapkan selama pelatihan untuk meningkatkan variasi data latih dan
mengurangi \textit{overfitting}. Teknik augmentasi yang digunakan meliputi:
\begin{itemize}
    \item \textit{MixUp} \citep{Zhang2018} ($\alpha = 0{,}2$, probabilitas 0,5 per \textit{batch}):
    menggabungkan dua sampel dan labelnya secara linear.
    \item \textit{CutMix} \citep{Yun2019} ($\alpha = 1{,}0$, probabilitas 0,5 per \textit{batch}):
    menempelkan potongan gambar dari satu sampel ke sampel lain dengan label proporsional.
    \item \textit{Label smoothing} 0,1: meratakan distribusi target agar model tidak terlalu
    percaya diri (\textit{overconfident}) pada satu kelas.
    \item \textit{Stochastic depth} (\textit{DropPath}) \citep{Huang2016} dengan jadwal
    linear $0 \rightarrow 0{,}2$: menonaktifkan sebagian jalur residual secara acak
    selama pelatihan sebagai bentuk regulasi.
    \item \textit{Weight EMA} (\textit{Exponential Moving Average}) dengan peluruhan 0,999:
    merata-ratakan bobot model secara eksponensial untuk mendapatkan model yang lebih stabil.
    \item \textit{Flip test-time augmentation} pada saat inferensi OOF: prediksi dirata-ratakan
    dari citra asli dan cerminan horizontalnya.
\end{itemize}

\subsection{Protokol Validasi Silang}

Protokol utama yang digunakan adalah \textit{StratifiedGroupKFold} dengan 5 \textit{fold}.
Pembagian ini memastikan bahwa: (1) distribusi kelas pada setiap \textit{fold} kira-kira
proporsional, dan (2) seluruh citra dari kelompok yang sama (kelompok citra serupa yang
terdeteksi melalui NCC) selalu berada di sisi yang sama. Prediksi \textit{out-of-fold} (OOF)
dari kelima fold digabungkan untuk menghitung metrik evaluasi akhir.

%--------------------------------------------------------------
\section{Metrik Evaluasi}
%--------------------------------------------------------------

\subsection{Metrik Klasifikasi}

Metrik utama yang digunakan adalah \textit{macro F1-score}, yaitu rata-rata dari nilai
F1-score masing-masing kelas tanpa memperhitungkan jumlah sampel per kelas. Metrik ini dipilih
karena ketidakseimbangan kelas pada dataset membuat akurasi keseluruhan tidak cukup informatif;
kelas dengan sampel lebih banyak dapat mendominasi akurasi meskipun kelas minoritas
berkinerja buruk. Formula \textit{macro F1-score} adalah:

\begin{equation}
  \text{Macro F1} = \frac{1}{K} \sum_{k=1}^{K} F1_k = \frac{1}{K} \sum_{k=1}^{K}
  \frac{2 \cdot P_k \cdot R_k}{P_k + R_k}
  \label{eq:macro_f1}
\end{equation}

\noindent di mana $K$ adalah jumlah kelas, $P_k$ adalah presisi, dan $R_k$ adalah \textit{recall}
untuk kelas ke-$k$.

Selain \textit{macro F1-score}, metrik tambahan yang dilaporkan mencakup akurasi keseluruhan,
\textit{macro precision}, \textit{macro recall}, dan F1-score per kelas.

\subsection{Metrik Segmentasi}

Kualitas segmentasi pembuluh darah dievaluasi menggunakan \textit{Dice coefficient}
\citep{Sorensen1948}, yang mengukur tumpang tindih antara masker prediksi dan masker
\textit{ground truth} dari anotasi pakar. Formula \textit{Dice coefficient} adalah:

\begin{equation}
  \text{Dice} = \frac{2 \times |A \cap B|}{|A| + |B|}
  \label{eq:dice}
\end{equation}

\noindent di mana $A$ adalah masker prediksi dan $B$ adalah masker \textit{ground truth}.
Evaluasi dilakukan pada 40 citra set pengujian yang tidak pernah terlihat selama optimasi
pipeline.
